\documentclass[11pt,a4paper]{article}

\usepackage[utf8]{inputenc}
\usepackage[indonesian]{babel}
\usepackage{amsmath,amssymb,amsfonts,amsthm}
\usepackage{geometry}
\usepackage{tcolorbox}
\usepackage{booktabs}
\usepackage{enumitem}
\usepackage{hyperref}
\usepackage{tikz}
\usepackage{pgfplots}
\pgfplotsset{compat=1.18}

\geometry{
    top=2.5cm,
    bottom=2.5cm,
    left=2.5cm,
    right=2.5cm
}

\tcbuselibrary{skins,breakable}

% Colors (sama persis dengan 06)
\definecolor{primaryblue}{RGB}{24, 76, 120}
\definecolor{accentblue}{RGB}{70, 130, 180}
\definecolor{boxbg}{RGB}{245, 248, 252}
\definecolor{darkgreen}{RGB}{34, 139, 34}
\definecolor{darkred}{RGB}{178, 34, 34}
\definecolor{darkorange}{RGB}{205, 92, 0}

% Custom theorem/box styles
\newtcolorbox{definitionbox}[1][]{
    enhanced,
    breakable,
    colback=boxbg,
    colframe=primaryblue,
    fonttitle=\bfseries\sffamily,
    title={#1},
    boxrule=0.8pt,
    titlerule=0pt,
    coltitle=white,
    colbacktitle=primaryblue,
    arc=3mm,
    boxsep=2mm
}

\newtcolorbox{examplebox}[1][]{
    enhanced,
    breakable,
    colback=white,
    colframe=accentblue,
    fonttitle=\bfseries\sffamily,
    title={#1},
    boxrule=0.8pt,
    arc=2mm,
    coltitle=white,
    colbacktitle=accentblue,
    boxsep=2mm
}

\newtcolorbox{notebox}[1][]{
    enhanced,
    breakable,
    colback=white,
    colframe=gray!60,
    fonttitle=\bfseries\sffamily,
    title={#1},
    boxrule=0.5pt,
    arc=2mm,
    coltitle=black!80,
    colbacktitle=gray!20,
    boxsep=2mm
}

\hypersetup{
    colorlinks=true,
    linkcolor=primaryblue,
    urlcolor=accentblue,
    citecolor=primaryblue
}

\title{\textbf{\LARGE Catatan Kuliah Teori Suku Bunga}\\[0.3em]
\Large Pertemuan 08: Anuitas yang Lebih Umum (Bagian 2)\\
\large (\textit{Continuous Annuities, Varying Annuities: Arithmetic \& Geometric, More General Varying Annuities, and Continuous Varying Annuities})}
\author{\textbf{Referensi:} Stephen G. Kellison, \textit{The Theory of Interest} (3rd ed., Bab 4) \& Slide Perkuliahan ITB}
\date{\today}

\begin{document}

\maketitle
\tableofcontents
\vspace{1em}
\noindent\rule{\textwidth}{0.4pt}
\vspace{1.5em}

\section{Pendahuluan \& Ruang Lingkup Materi}

Pada Bab 05 dan 06, kita telah mempelajari anuitas dasar di mana pembayaran dilakukan secara diskrit dengan jumlah pembayaran yang konstan (\textit{level annuity}). Kemudian pada Bab 07, kita memperluas analisis ke kasus di mana frekuensi pembayaran berbeda dari frekuensi pemajemukan suku bunga, namun pembayaran di setiap periode tetap sama.

Dalam praktik industri keuangan, perbankan, dan aktuarial modern, terdapat dua kondisi lanjutan yang sering dijumpai:
\begin{enumerate}[label=\textbf{\arabic*.}]
    \item \textbf{Anuitas Kontinu (\textit{Continuous Annuities}):} Pembayaran angsuran atau penerimaan dana dilakukan secara terus-menerus sepanjang waktu (frekuensi pembayaran menuju tak hingga, $m \to \infty$). Hal ini sangat penting secara teoritis dan analitis, serta menjadi aproksimasi yang sangat baik untuk transaksi harian, klaim/premi asuransi jiwa kontinu, dan arus kas operasional korporasi.
    \item \textbf{Anuitas Bervariasi (\textit{Non-Level / Varying Annuities}):} Pembayaran tidak tetap melainkan berubah dari waktu ke waktu:
    \begin{itemize}
        \item \textbf{Pola Aritmatika:} Pembayaran bertambah atau berkurang dengan nominal tetap ($Q$) setiap periode (anuitas meningkat $(Ia)$ dan menurun $(Da)$).
        \item \textbf{Pola Geometri:} Pembayaran bertumbuh atau menyusut dengan persentase majemuk ($k$) setiap periode (penyesuaian inflasi COLA dan valuasi dividen bertumbuh).
        \item \textbf{Variasi Lebih Umum:} Variasi pembayaran dengan frekuensi pembayaran yang berbeda dari frekuensi konversi bunga ($k > 1$ atau $m > 1$).
        \item \textbf{Variasi Kontinu:} Pembayaran dilakukan secara kontinu dengan laju yang bervariasi mengikuti fungsi waktu $f(t)$ atau dengan \textit{force of interest} $\delta(t)$ yang dinamis.
    \end{itemize}
\end{enumerate}

---

\section{Anuitas Kontinu Sederhana (\textit{Continuous Annuities})}

\subsection{Definisi \& Konsep Dasar}

\begin{definitionbox}[Anuitas Kontinu Berjangka $n$ Periode]
Suatu \textbf{anuitas kontinu} adalah anuitas di mana pembayaran dilakukan secara terus-menerus (\textit{continuously}) sepanjang waktu, sedemikian rupa sehingga jumlah total yang dibayarkan selama setiap satu periode konversi suku bunga tepat bernilai $1$.
\end{definitionbox}

Konsep ini diperoleh dari limit anuitas diskrit ketika frekuensi pembayaran $m$ menuju tak hingga ($m \to \infty$). Dalam interval waktu infinitesimal $[t, t+dt]$, besar pembayaran adalah $dt$.
Karena pembayaran sebesar $dt$ terjadi pada waktu eksak $t$, nilai sekarang (\textit{present value}) dari pembayaran tersebut adalah:
\begin{equation}
    d(PV) = v^t \, dt = e^{-\delta t} \, dt
\end{equation}
di mana $\delta = \ln(1+i) = -\ln v$ adalah \textit{force of interest} konstan.

\begin{center}
\begin{tikzpicture}[scale=1.0, >=stealth]
    \draw[->, thick] (-0.5,0) -- (10.5,0) node[right] {Waktu ($t$)};
    \draw[thick] (0,0.15) -- (0,-0.15) node[below] {$0$};
    \draw[thick] (2,0.1) -- (2,-0.1) node[below] {$1$};
    \draw[thick] (4,0.1) -- (4,-0.1) node[below] {$2$};
    \draw[dotted, thick] (4.5,0) -- (7.5,0);
    \draw[thick] (8,0.15) -- (8,-0.15) node[below] {$n$};

    % Arsiran aliran kontinu
    \fill[primaryblue!20] (0,0) rectangle (8,1.2);
    \draw[thick, primaryblue] (0,1.2) -- (8,1.2) node[midway, above] {\small Laju aliran kas konstan: $1$ per periode waktu};
    \draw[thick, primaryblue] (0,0) -- (0,1.2);
    \draw[thick, primaryblue] (8,0) -- (8,1.2);

    % Elemen infinitesimal dt
    \fill[darkred!40] (3.2,0) rectangle (3.5,1.2);
    \draw[darkred, thick] (3.35, 1.3) -- (3.35, 1.9) node[above] {\footnotesize Pembayaran $1 \cdot dt$ pada waktu $t$};
    \draw[->, darkred, thick] (3.35, 0) -- (0, -0.6) node[below] {\footnotesize Didiskonto ke $t=0$: $v^t dt$};
\end{tikzpicture}
\end{center}

---

\subsection{Penurunan Nilai Sekarang (\texorpdfstring{$\bar{a}_{\overline{n}|}$}{a_bar_n})}

Simbol nilai sekarang dari anuitas kontinu berjangka $n$ periode adalah $\mathbf{\bar{a}_{\overline{n}|}}$:
\begin{equation}
    \bar{a}_{\overline{n}|} = \int_0^n v^t \, dt = \int_0^n e^{t \ln v} \, dt = \left[ \frac{v^t}{\ln v} \right]_0^n = \frac{v^n - 1}{\ln v}
\end{equation}
Mengingat definisi \textit{force of interest}:
\begin{equation}
    \delta = \frac{a'(t)}{a(t)} = \frac{(1+i)^t \ln(1+i)}{(1+i)^t} = \ln(1+i) = -\ln\left(\frac{1}{1+i}\right) = -\ln v
\end{equation}
Maka $\ln v = -\delta$, sehingga diperoleh rumus utama:
\begin{equation}
    \mathbf{\bar{a}_{\overline{n}|} = \frac{1 - v^n}{\delta}}
\end{equation}

\subsubsection{Konsistensi dengan Anuitas Diskrit}
Rumus ini menunjukkan kekonsistenan bentuk yang sangat harmonis dalam teori suku bunga:
\begin{align}
    \text{Anuitas Akhir Diskrit:} \quad a_{\overline{n}|} &= \frac{1 - v^n}{i} \\
    \text{Anuitas Kontinu:} \quad \mathbf{\bar{a}_{\overline{n}|}} &= \mathbf{\frac{1 - v^n}{\delta}} \\
    \text{Anuitas Awal Diskrit:} \quad \ddot{a}_{\overline{n}|} &= \frac{1 - v^n}{d}
\end{align}
Bentuk yang sama juga diperoleh dari limit anuitas diskrit saat $m \to \infty$:
\begin{align}
    \bar{a}_{\overline{n}|} &= \lim_{m \to \infty} a_{\overline{n}|}^{(m)} = \lim_{m \to \infty} \frac{1 - v^n}{i^{(m)}} = \frac{1 - v^n}{\delta} \\
    \bar{a}_{\overline{n}|} &= \lim_{m \to \infty} \ddot{a}_{\overline{n}|}^{(m)} = \lim_{m \to \infty} \frac{1 - v^n}{d^{(m)}} = \frac{1 - v^n}{\delta}
\end{align}

---

\subsection{Penurunan Nilai Akumulasi (\texorpdfstring{$\bar{s}_{\overline{n}|}$}{s_bar_n})}

Nilai akumulasi pada akhir periode $t = n$ dinotasikan dengan $\mathbf{\bar{s}_{\overline{n}|}}$:
\begin{equation}
    \bar{s}_{\overline{n}|} = \int_0^n (1+i)^t \, dt = \int_0^n e^{t \ln(1+i)} \, dt = \left[ \frac{(1+i)^t}{\ln(1+i)} \right]_0^n = \mathbf{\frac{(1+i)^n - 1}{\delta} = \frac{e^{n\delta} - 1}{\delta}}
\end{equation}
Dapat juga diperoleh dari hubungan akumulasi dasar:
\begin{equation}
    \bar{s}_{\overline{n}|} = \bar{a}_{\overline{n}|} \cdot (1+i)^n = \bar{a}_{\overline{n}|} \cdot e^{n\delta}
\end{equation}
dan berlaku sifat limit:
\begin{equation}
    \bar{s}_{\overline{n}|} = \lim_{m \to \infty} s_{\overline{n}|}^{(m)} = \lim_{m \to \infty} \ddot{s}_{\overline{n}|}^{(m)}
\end{equation}

---

\subsection{Interpretasi Dinamik \& Persamaan Diferensial (Slide Hal. 4--5)}

\subsubsection{Diferensiasi Terhadap Nilai Akumulasi}
Gunakan Teorema Dasar Kalkulus Pertama:
\begin{equation}
    \frac{d}{dt} \bar{s}_{\overline{t}|} = \frac{d}{dt} \int_0^t (1+i)^x \, dx = (1+i)^t = 1 + \delta \, \bar{s}_{\overline{t}|}
\end{equation}

\begin{notebox}[Interpretasi Finansial Dinamika Saldo Dana]
Misalkan suatu dana investasi di mana uang disimpan secara kontinu dengan laju $1$ per periode konversi suku bunga. Saldo saat $t$ adalah $\bar{s}_{\overline{t}|}$. Laju perubahan saldo dana terdiri atas:
\begin{enumerate}
    \item Setoran baru masuk dengan laju $1$ per periode.
    \item Bunga yang diperoleh dengan laju $\delta$ terhadap saldo yang ada ($\delta \bar{s}_{\overline{t}|}$).
\end{enumerate}
\end{notebox}

\subsubsection{Diferensiasi Terhadap Nilai Sekarang \& Solusi PD Linier}
Untuk nilai sekarang anuitas kontinu:
\begin{equation}
    \frac{d}{dt} \bar{a}_{\overline{t}|} = \frac{d}{dt} \int_0^t v^x \, dx = v^t = 1 - \delta \, \bar{a}_{\overline{t}|}
\end{equation}
Misalkan $y = \bar{a}_{\overline{t}|}$, maka persamaan menjadi:
\begin{equation}
    y' = 1 - \delta y \iff y' + \delta y = 1
\end{equation}
Ini adalah \textbf{Persamaan Diferensial Linier Orde Satu Tak Homogen}:
\begin{itemize}
    \item Persamaan homogen: $y' + \delta y = 0 \implies y_h = c e^{-\delta t}$.
    \item Solusi partikular: Misalkan $y_p = a$ (konstanta), maka $0 + \delta a = 1 \implies a = \frac{1}{\delta}$.
    \item Solusi umum: $y = y_h + y_p = c e^{-\delta t} + \frac{1}{\delta}$.
    \item Syarat awal: Untuk $t = 0 \implies y(0) = \bar{a}_{\overline{0}|} = 0$:
    \begin{equation*}
        0 = c e^0 + \frac{1}{\delta} \implies c = -\frac{1}{\delta}
    \end{equation*}
\end{itemize}
Maka solusi tunggalnya adalah:
\begin{equation}
    y = \bar{a}_{\overline{t}|} = -\frac{1}{\delta} e^{-\delta t} + \frac{1}{\delta} = \mathbf{\frac{1 - e^{-\delta t}}{\delta} = \frac{1 - v^t}{\delta}}
\end{equation}

---

\subsection{Contoh Soal 1: Penentuan Tingkat Bunga Kontinu (Slide Hal. 6)}

\begin{examplebox}[Contoh 1 (Slide Perkuliahan Hal. 6)]
Hitung \textit{force of interest} $\delta$ dari persamaan:
\begin{equation*}
    \bar{s}_{\overline{20}|} = 3 \, \bar{s}_{\overline{10}|}
\end{equation*}
\end{examplebox}

\begin{proof}[\textbf{Penyelesaian:}]
Substitusikan rumus $\bar{s}_{\overline{n}|} = \dfrac{e^{n\delta} - 1}{\delta}$:
\begin{equation*}
    \frac{e^{20\delta} - 1}{\delta} = 3 \left( \frac{e^{10\delta} - 1}{\delta} \right)
\end{equation*}
Kalikan kedua ruas dengan $\delta$:
\begin{equation*}
    e^{20\delta} - 1 = 3 e^{10\delta} - 3 \implies e^{20\delta} - 3 e^{10\delta} + 2 = 0
\end{equation*}
Faktorkan persamaan kuadrat terhadap $e^{10\delta}$:
\begin{equation*}
    (e^{10\delta} - 2)(e^{10\delta} - 1) = 0
\end{equation*}
\begin{itemize}
    \item Jika $e^{10\delta} - 1 = 0 \implies e^{10\delta} = 1 \implies \delta = 0$. Solusi ini ditolak (suku bunga nol tidak berlaku dalam transaksi keuangan).
    \item Jika $e^{10\delta} - 2 = 0 \implies e^{10\delta} = 2$.
    Ambil logaritma natural:
    \begin{equation*}
        10\delta = \ln 2 \implies \delta = \frac{\ln 2}{10} \approx \frac{0.693147}{10} = \mathbf{0.0693} = \mathbf{6.93\%}
    \end{equation*}
\end{itemize}
Jadi, \textit{force of interest} yang memenuhi adalah $\mathbf{\delta \approx 6.93\%}$.
\end{proof}

---

\section{Anuitas Bervariasi secara Aritmatika (\textit{Arithmetic Varying Annuities})}

\subsection{Definisi \& Penurunan Rumus Umum (\texorpdfstring{$P, Q$}{P, Q})}

\begin{definitionbox}[Anuitas Bervariasi secara Aritmatika]
Misalkan suatu annuity-immediate berjangka $n$ periode dengan suku bunga $i$ per periode:
\begin{itemize}
    \item Pembayaran pertama bernilai $P$ ($P > 0$).
    \item Pembayaran bertambah (atau berkurang) sebesar $Q$ di setiap periode berikutnya.
    \item Barisan pembayaran pada $t = 1, 2, \dots, n$ adalah:
    \begin{equation*}
        P, \quad P+Q, \quad P+2Q, \quad \dots, \quad P+(n-1)Q
    \end{equation*}
    dengan syarat kelayakan finansial $P + (n-1)Q > 0$.
\end{itemize}
\end{definitionbox}

Misalkan $A$ adalah nilai sekarang:
\begin{equation}
    A = P v + (P+Q) v^2 + (P+2Q) v^3 + \dots + [P + (n-1)Q] v^n \label{eq:aritA}
\end{equation}
Kalikan kedua ruas dengan $(1+i)$:
\begin{equation}
    (1+i)A = P + (P+Q) v + (P+2Q) v^2 + \dots + [P + (n-1)Q] v^{n-1} \label{eq:aritB}
\end{equation}
Kurangkan Persamaan~(\ref{eq:aritA}) dari Persamaan~(\ref{eq:aritB}):
\begin{align*}
    i A &= P + Q v + Q v^2 + \dots + Q v^{n-1} - [P + (n-1)Q] v^n \\
        &= P(1 - v^n) + Q(v + v^2 + \dots + v^{n-1}) - (n-1)Q v^n \\
        &= P(1 - v^n) + Q(v + v^2 + \dots + v^n) - Q n v^n \\
        &= P(1 - v^n) + Q \cdot a_{\overline{n}|} - Q n v^n
\end{align*}
Bagi kedua ruas dengan suku bunga $i$:
\begin{equation}
    \mathbf{A = P \left( \frac{1 - v^n}{i} \right) + Q \left( \frac{a_{\overline{n}|} - n v^n}{i} \right) = P \, a_{\overline{n}|} + Q \left( \frac{a_{\overline{n}|} - n v^n}{i} \right)}
\end{equation}
Nilai akumulasi ($AV$) pada akhir periode $n$:
\begin{equation}
    \mathbf{AV = A(1+i)^n = P \, s_{\overline{n}|} + Q \left( \frac{s_{\overline{n}|} - n}{i} \right)}
\end{equation}

---

\subsection{Kasus Khusus 1: Anuitas Meningkat Standar (\texorpdfstring{$(Ia)_{\overline{n}|}$}{(Ia)_n})}

Jika pembayaran bernilai $1, 2, 3, \dots, n$ pada akhir setiap periode (di mana $P = 1, Q = 1$):
\begin{equation}
    (Ia)_{\overline{n}|} = a_{\overline{n}|} + \frac{a_{\overline{n}|} - n v^n}{i} = \frac{i \, a_{\overline{n}|} + a_{\overline{n}|} - n v^n}{i}
\end{equation}
Karena $i \, a_{\overline{n}|} = 1 - v^n$, pembilang menjadi:
\begin{equation*}
    (1 - v^n) + a_{\overline{n}|} - n v^n = (1 + a_{\overline{n}|} - v^n) - n v^n = \ddot{a}_{\overline{n}|} - n v^n
\end{equation*}
Sehingga diperoleh rumus standar:
\begin{equation}
    \mathbf{(Ia)_{\overline{n}|} = \frac{\ddot{a}_{\overline{n}|} - n v^n}{i}}
\end{equation}
Nilai akumulasi pada akhir periode $n$:
\begin{equation}
    \mathbf{(Is)_{\overline{n}|} = (Ia)_{\overline{n}|} (1+i)^n = \frac{\ddot{s}_{\overline{n}|} - n}{i}}
\end{equation}

\begin{center}
\begin{tikzpicture}[scale=1.0, >=stealth]
    \draw[->, thick] (-0.5,0) -- (10.5,0) node[right] {Waktu ($t$)};
    \foreach \t in {0,1,2,3,4} {
        \draw[thick] (\t*1.8, 0.1) -- (\t*1.8, -0.1) node[below] {\small $\t$};
    }
    \draw[dotted, thick] (7.6,0) -- (8.8,0);
    \draw[thick] (9.5, 0.1) -- (9.5, -0.1) node[below] {\small $n$};

    % Batang arus kas bertingkat
    \draw[->, very thick, primaryblue] (1.8, 0.1) -- (1.8, 0.7) node[above] {\small $\$1$};
    \draw[->, very thick, primaryblue] (3.6, 0.1) -- (3.6, 1.3) node[above] {\small $\$2$};
    \draw[->, very thick, primaryblue] (5.4, 0.1) -- (5.4, 1.9) node[above] {\small $\$3$};
    \draw[->, very thick, primaryblue] (7.2, 0.1) -- (7.2, 2.5) node[above] {\small $\$4$};
    \draw[->, very thick, primaryblue] (9.5, 0.1) -- (9.5, 3.2) node[above] {\small $\$n$};
\end{tikzpicture}
\end{center}

---

\subsection{Kasus Khusus 2: Anuitas Menurun Standar (\texorpdfstring{$(Da)_{\overline{n}|}$}{(Da)_n})}

Jika pembayaran bernilai $n, n-1, n-2, \dots, 2, 1$ pada akhir setiap periode (di mana $P = n, Q = -1$):
\begin{align*}
    (Da)_{\overline{n}|} &= n \, a_{\overline{n}|} - \frac{a_{\overline{n}|} - n v^n}{i} = \frac{n i a_{\overline{n}|} - a_{\overline{n}|} + n v^n}{i} \\
                        &= \frac{n(1 - v^n) - a_{\overline{n}|} + n v^n}{i} = \frac{n - n v^n - a_{\overline{n}|} + n v^n}{i}
\end{align*}
Suku $n v^n$ saling menghilangkan, menghasilkan formula sederhana:
\begin{equation}
    \mathbf{(Da)_{\overline{n}|} = \frac{n - a_{\overline{n}|}}{i}}
\end{equation}
Nilai akumulasi pada akhir periode $n$:
\begin{equation}
    \mathbf{(Ds)_{\overline{n}|} = (Da)_{\overline{n}|} (1+i)^n = \frac{n(1+i)^n - s_{\overline{n}|}}{i}}
\end{equation}

\begin{notebox}[Aturan Anuitas di Muka (\textit{Annuity-Due}) \& Teorema Komplementer]
\begin{enumerate}
    \item Untuk \textbf{annuity-due}, penyebut $i$ diganti dengan tingkat diskon $d$:
    \begin{equation}
        (I\ddot{a})_{\overline{n}|} = \frac{\ddot{a}_{\overline{n}|} - n v^n}{d}, \qquad (D\ddot{a})_{\overline{n}|} = \frac{n - a_{\overline{n}|}}{d}
    \end{equation}
    \item \textbf{Relasi Komplementer:} Karena pada setiap periode $t$, $t + (n - t + 1) = n + 1$, maka:
    \begin{equation}
        \mathbf{(Ia)_{\overline{n}|} + (Da)_{\overline{n}|} = (n+1) \, a_{\overline{n}|}}
    \end{equation}
\end{enumerate}
\end{notebox}

---

\subsection{Perpetuitas Bervariasi secara Aritmatika (Slide Hal. 11)}

Dengan mengambil limit $n \to \infty$ dari rumus umum $(P, Q)$:
\begin{equation}
    \lim_{n \to \infty} \left[ P a_{\overline{n}|} + Q \left( \frac{a_{\overline{n}|} - n v^n}{i} \right) \right] = P \left(\frac{1}{i}\right) + Q \left(\frac{1/i - 0}{i}\right) = \mathbf{\frac{P}{i} + \frac{Q}{i^2}}
\end{equation}
karena $\lim\limits_{n \to \infty} a_{\overline{n}|} = \dfrac{1}{i}$ dan $\lim\limits_{n \to \infty} n v^n = \lim\limits_{n \to \infty} \dfrac{n}{(1+i)^n} = 0$.

Untuk perpetuitas meningkat standar ($P=1, Q=1$):
\begin{equation}
    \mathbf{(Ia)_{\overline{\infty}|} = \frac{1}{i} + \frac{1}{i^2} = \frac{1}{i \, d}}
\end{equation}

---

\section{Anuitas Bervariasi secara Geometrik (\textit{Geometric Varying Annuities})}

\subsection{Definisi \& Penurunan Nilai Sekarang}

\begin{definitionbox}[Anuitas Geometrik Berjangka $n$ Periode]
Suatu anuitas di mana pembayaran pertama adalah $1$, dan setiap pembayaran berikutnya membesar secara geometri dengan rasio $1+k$ ($k > -1$).
Aliran kas: $1, (1+k), (1+k)^2, \dots, (1+k)^{n-1}$.
\end{definitionbox}

Nilai sekarang anuitas akhir adalah deret geometri berjangka $n$:
\begin{align*}
    PV &= v + (1+k) v^2 + (1+k)^2 v^3 + \dots + (1+k)^{n-1} v^n \\
       &= v \left[ 1 + (1+k)v + \dots + (1+k)^{n-1} v^{n-1} \right] \\
       &= v \left[ \frac{1 - \left(\frac{1+k}{1+i}\right)^n}{1 - \frac{1+k}{1+i}} \right] = \frac{1}{1+i} \left[ \frac{1 - \left(\frac{1+k}{1+i}\right)^n}{\frac{(1+i) - (1+k)}{1+i}} \right]
\end{align*}
Sehingga:
\begin{equation}
    \mathbf{PV = \frac{1 - \left(\frac{1+k}{1+i}\right)^n}{i - k}} \qquad (k \ne i)
\end{equation}
\begin{notebox}[Kasus Khusus $k = i$]
Jika $k = i$, maka pembagian $i - k$ bernilai nol. Namun dari deret awal:
\begin{equation*}
    PV = v + v + \dots + v = n \, v = \mathbf{\frac{n}{1+i}}
\end{equation*}
\end{notebox}

Nilai akumulasi pada akhir periode $n$:
\begin{equation}
    \mathbf{AV = PV \cdot (1+i)^n = \frac{(1+i)^n - (1+k)^n}{i - k}} \qquad (k \ne i)
\end{equation}

---

\subsection{Metode Suku Bunga Pengganti Ekuivalen (\texorpdfstring{$\tilde{i}$}{i_tilde}) (Slide Hal. 13)}

Jika $k < i$, definisikan suku bunga baru $\tilde{i}$ sedemikian rupa sehingga:
\begin{equation}
    1 + \tilde{i} = \frac{1+i}{1+k} \iff \tilde{i} = \frac{1+i}{1+k} - 1 = \mathbf{\frac{i - k}{1 + k}}
\end{equation}
Maka untuk \textbf{annuity-due} dengan rasio bertumbuh:
\begin{equation}
    1 + \frac{1+k}{1+i} + \dots + \left(\frac{1+k}{1+i}\right)^{n-1} = 1 + \frac{1}{1+\tilde{i}} + \dots + \frac{1}{(1+\tilde{i})^{n-1}} = \mathbf{\ddot{a}_{\overline{n}|\tilde{i}}}
\end{equation}
Artinya: Nilai sekarang dari anuitas $n$ periode dengan pembayaran membesar secara geometrik rasio $1+k$ ($k < i$) adalah \textbf{sama dengan nilai sekarang anuitas pembayaran tetap yang dievaluasi pada suku bunga $\tilde{i}$}.

Untuk \textbf{perpetuitas geometrik} ($n \to \infty$ dan $k < i$):
\begin{equation}
    \mathbf{PV_{\text{perp-geometrik}} = \frac{1}{i - k}}
\end{equation}

---

\subsection{Contoh Soal 2: Valuasi Anuitas Bertumbuh Geometrik (Slide Hal. 15)}

\begin{examplebox}[Contoh 2 (Slide Perkuliahan Hal. 15)]
Suatu anuitas memiliki pembayaran per tahun sebanyak 20 kali. Pembayaran pertama adalah $\$1.000$. Pembayaran meningkat sedemikian rupa sehingga setiap pembayaran $4\%$ lebih besar daripada pembayaran sebelumnya. Hitung nilai sekarang dari anuitas ini jika suku bunga adalah $7\%$.
\end{examplebox}

\begin{proof}[\textbf{Penyelesaian:}]
Data: $n = 20$, $R_1 = 1000$, $k = 0.04$, $i = 0.07$.
Nilai sekarang anuitasnya adalah:
\begin{align*}
    PV &= 1000 \left[ \frac{1}{1.07} + \frac{1.04}{(1.07)^2} + \dots + \frac{(1.04)^{19}}{(1.07)^{20}} \right] \\
       &= 1000 \left[ \frac{1 - \left(\frac{1.04}{1.07}\right)^{20}}{0.07 - 0.04} \right] \\
       &= 1000 \left[ \frac{1 - (0.9719626)^{20}}{0.03} \right] \\
       &= 1000 \left[ \frac{1 - 0.566228}{0.03} \right] = 1000 \left[ \frac{0.433772}{0.03} \right] = 1000 \times 14.45907 = \mathbf{\$14.459,07}
\end{align*}
Jadi, nilai sekarang anuitas tersebut adalah \textbf{\$14.459} (sesuai slide).
\end{proof}

---

\section{Anuitas Bervariasi yang Lebih Umum (\textit{More General Varying Annuities})}

\subsection{Frekuensi Pembayaran Lebih Sedikit daripada Konversi Bunga (\texorpdfstring{$k > 1$}{k > 1})}

Misalkan:
\begin{itemize}
    \item $k$: Banyaknya periode konversi bunga dalam 1 periode pembayaran ($k \in \mathbb{Z}^+$).
    \item $n$: Total durasi anuitas diukur dalam periode konversi bunga.
    \item Frekuensi pembayaran adalah $m = \dfrac{n}{k} \in \mathbb{Z}^+$.
    \item Pembayaran dilakukan sebesar $1$ di $t=k$, $2$ di $t=2k$, ..., $n/k$ di $t=n$.
\end{itemize}

\subsubsection{Penurunan Matematis Nilai Sekarang (Slide Hal. 17)}
Nilai sekarang $A$:
\begin{equation}
    A = v^k + 2 v^{2k} + \dots + \left(\frac{n}{k}-1\right) v^{n-k} + \left(\frac{n}{k}\right) v^n \label{eq:gen1}
\end{equation}
Kalikan dengan $(1+i)^k$:
\begin{equation}
    (1+i)^k A = 1 + 2 v^k + \dots + \left(\frac{n}{k}\right) v^{n-k} \label{eq:gen2}
\end{equation}
Kurangkan Persamaan~(\ref{eq:gen1}) dari Persamaan~(\ref{eq:gen2}):
\begin{equation*}
    A \left[ (1+i)^k - 1 \right] = 1 + v^k + v^{2k} + \dots + v^{n-k} - \frac{n}{k} v^n
\end{equation*}
Perhatikan bahwa $1 + v^k + \dots + v^{n-k} = \dfrac{a_{\overline{n}|}}{a_{\overline{k}|}}$, dan $(1+i)^k - 1 = i \cdot s_{\overline{k}|}$.
Maka:
\begin{equation}
    \mathbf{A = \frac{\dfrac{a_{\overline{n}|}}{a_{\overline{k}|}} - \dfrac{n}{k} v^n}{i \, s_{\overline{k}|}}}
\end{equation}
Persamaan ini merupakan versi umum dari $(Ia)_{\overline{n}|} = \dfrac{\ddot{a}_{\overline{n}|} - n v^n}{i}$.

---

\subsection{Frekuensi Pembayaran Lebih Sering daripada Konversi Bunga (\texorpdfstring{$m > 1$}{m > 1})}

Ketika pembayaran dilakukan $m$ kali per periode bunga, terdapat 2 kasus penting:

\subsubsection{Kasus A: Pembayaran Konstan Tiap Periode, Kenaikan Sekali per Periode (Slide Hal. 18)}
Pembayaran dilakukan sebesar $\frac{t}{m}$ di setiap $1/m$ periode pada tahun ke-$t$ (sehingga total pembayaran tahunan adalah $t$):
\begin{equation}
    \mathbf{(Ia)_{\overline{n}|}^{(m)} = \frac{\ddot{a}_{\overline{n}|} - n v^n}{i^{(m)}}}
\end{equation}

\subsubsection{Kasus B: Laju Pembayaran Meningkat pada Setiap Sub-Periode (Slide Hal. 19)}
Pembayaran pada waktu $\frac{t}{m}$ bernilai $\frac{t}{m^2}$ (pembayaran pertama $\frac{1}{m^2}$, kedua $\frac{2}{m^2}$, dan seterusnya):
\begin{equation}
    (I^{(m)} a)_{\overline{n}|}^{(m)} = \frac{1}{m^2} v^{1/m} + \frac{2}{m^2} v^{2/m} + \dots + \frac{nm}{m^2} v^{nm/m}
\end{equation}
\begin{equation}
    \mathbf{(I^{(m)} a)_{\overline{n}|}^{(m)} = \frac{\ddot{a}_{\overline{n}|}^{(m)} - n v^n}{i^{(m)}}}
\end{equation}

---

\subsection{Contoh Soal 3: Perpetuitas Meningkat Multi-Tahun (Slide Hal. 20)}

\begin{examplebox}[Contoh 3 (Slide Perkuliahan Hal. 20)]
Cari nilai sekarang dari suatu perpetuitas yang membayar:
\begin{itemize}
    \item $\$1$ pada akhir tahun ke-3,
    \item $\$2$ pada akhir tahun ke-6,
    \item $\$3$ pada akhir tahun ke-9, dan seterusnya.
\end{itemize}
\end{examplebox}

\begin{proof}[\textbf{Penyelesaian:}]
Nilai sekarang perpetuitas tersebut adalah:
\begin{equation*}
    A = v^3 + 2 v^6 + 3 v^9 + \dots
\end{equation*}
Kalikan kedua ruas dengan $v^3$:
\begin{equation*}
    v^3 A = v^6 + 2 v^9 + 3 v^{12} + \dots
\end{equation*}
Kurangkan persamaan kedua dari yang pertama:
\begin{equation*}
    (1 - v^3) A = v^3 + v^6 + v^9 + \dots
\end{equation*}
Ruas kanan adalah deret geometri tak hingga:
\begin{equation*}
    (1 - v^3) A = \frac{v^3}{1 - v^3}
\end{equation*}
Sehingga:
\begin{equation*}
    \mathbf{A = \frac{v^3}{(1 - v^3)^2}}
\end{equation*}
Secara alternatif, jika kita mendefinisikan suku bunga 3-tahunan $j = (1+i)^3 - 1$, maka $v_j = v^3$ dan $A = (Ia)_{\overline{\infty}|j} = \frac{1}{j} + \frac{1}{j^2} = \frac{1+j}{j^2} = \frac{v_j}{(1 - v_j)^2} = \frac{v^3}{(1 - v^3)^2}$.
\end{proof}

---

\section{Anuitas Bervariasi secara Kontinu (\textit{Continuous Varying Annuities})}

\subsection{Anuitas Meningkat Kontinu Standar (\texorpdfstring{$(\bar{I}\bar{a})_{\overline{n}|}$}{(I_bar a_bar)_n})}

Misalkan anuitas dibayar kontinu selama $n$ periode konversi bunga dengan laju pembayaran pada waktu $t$ adalah $t$ ($f(t) = t$).
Nilai sekarangnya adalah:
\begin{equation}
    (\bar{I}\bar{a})_{\overline{n}|} = \int_0^n t \, v^t \, dt = \mathbf{\frac{\bar{a}_{\overline{n}|} - n v^n}{\delta}}
\end{equation}
Rumus ini diperoleh dari limit anuitas diskrit:
\begin{equation}
    (\bar{I}\bar{a})_{\overline{n}|} = \lim_{m \to \infty} (I^{(m)} a)_{\overline{n}|}^{(m)} = \lim_{m \to \infty} \frac{\ddot{a}_{\overline{n}|}^{(m)} - n v^n}{i^{(m)}} = \frac{\bar{a}_{\overline{n}|} - n v^n}{\delta}
\end{equation}
Nilai akumulasinya adalah:
\begin{equation}
    (\bar{I}\bar{s})_{\overline{n}|} = (\bar{I}\bar{a})_{\overline{n}|} \cdot (1+i)^n = \mathbf{\frac{\bar{s}_{\overline{n}|} - n}{\delta}}
\end{equation}

Untuk anuitas menurun kontinu standar ($f(t) = n - t$):
\begin{equation}
    \mathbf{(\bar{D}\bar{a})_{\overline{n}|} = \int_0^n (n - t) v^t \, dt = \frac{n - \bar{a}_{\overline{n}|}}{\delta}}
\end{equation}

---

\subsection{Laju Pembayaran Kontinu Umum \texorpdfstring{$f(t)$}{f(t)} \& \texorpdfstring{$\delta(t)$}{delta(t)}}

\begin{definitionbox}[Present Value Aliran Pembayaran Kontinu Umum]
\begin{enumerate}
    \item Jika pembayaran pada waktu $t$ dilakukan dengan laju $f(t) dt$, maka nilai sekarang $n$-periode adalah:
    \begin{equation}
        \mathbf{PV = \int_0^n f(t) \, v^t \, dt = \int_0^n f(t) \, e^{-\delta t} \, dt}
    \end{equation}
    \item Jika suku bunga juga bervariasi dengan \textit{force of interest} $\delta(r)$, maka:
    \begin{equation}
        \mathbf{PV = \int_0^n f(t) \exp\left( -\int_0^t \delta(r) \, dr \right) dt}
    \end{equation}
\end{enumerate}
\end{definitionbox}

---

\subsection{Contoh Soal 4: Anuitas Kontinu dengan Laju Kuadratik \texorpdfstring{$f(t) = t^2$}{f(t) = t^2} (Slide Hal. 22)}

\begin{examplebox}[Contoh 4 (Slide Perkuliahan Hal. 22)]
Tentukan \textit{present value} dari \textit{continuously increasing annuity} dengan term $n$ tahun, jika \textit{force of interest} adalah $\delta$ dan rate pembayaran saat waktu $t$ adalah $t^2$ per tahun.
\end{examplebox}

\begin{proof}[\textbf{Penyelesaian Lengkap via Integrasi Parsial:}]
Integral nilai sekarang adalah:
\begin{equation*}
    PV = \int_0^n t^2 e^{-\delta t} \, dt
\end{equation*}
Gunakan integrasi parsial berulang:
\begin{itemize}
    \item Integrasi ke-1: Misal $u = t^2 \implies du = 2t dt$, $dv = e^{-\delta t} dt \implies v = -\frac{1}{\delta} e^{-\delta t}$:
    \begin{equation*}
        \int t^2 e^{-\delta t} dt = -\frac{t^2}{\delta} e^{-\delta t} + \frac{2}{\delta} \int t e^{-\delta t} dt
    \end{equation*}
    \item Integrasi ke-2: Dari rumus anuitas meningkat $(\bar{I}\bar{a})$: $\int t e^{-\delta t} dt = -\frac{t}{\delta} e^{-\delta t} - \frac{1}{\delta^2} e^{-\delta t}$.
\end{itemize}
Maka antiturunannya adalah:
\begin{equation*}
    \int t^2 e^{-\delta t} dt = -e^{-\delta t} \left[ \frac{t^2}{\delta} + \frac{2t}{\delta^2} + \frac{2}{\delta^3} \right]
\end{equation*}
Evaluasi pada batas $[0, n]$:
\begin{align*}
    PV &= \left( -e^{-\delta n} \left[ \frac{n^2}{\delta} + \frac{2n}{\delta^2} + \frac{2}{\delta^3} \right] \right) - \left( -e^0 \left[ 0 + 0 + \frac{2}{\delta^3} \right] \right) \\
       &= \mathbf{\frac{2}{\delta^3} - e^{-n\delta} \left[ \frac{n^2}{\delta} + \frac{2n}{\delta^2} + \frac{2}{\delta^3} \right]}
\end{align*}
Dalam simbol aktuaria dapat disederhanakan menjadi:
\begin{equation*}
    PV = \frac{2(1 - v^n)}{\delta^3} - \frac{2n v^n}{\delta^2} - \frac{n^2 v^n}{\delta} = \frac{2\bar{a}_{\overline{n}|} - 2n v^n}{\delta^2} - \frac{n^2 v^n}{\delta} = \mathbf{\frac{2(\bar{I}\bar{a})_{\overline{n}|} - n^2 v^n}{\delta}}
\end{equation*}
\end{proof}

---

\section{Tabel Ringkasan Rumus Anuitas Bab 08}

\begin{center}
\small
\resizebox{\textwidth}{!}{
\begin{tabular}{lcccc}
\toprule
\textbf{Jenis Anuitas} & \textbf{Aliran Kas pada Waktu $t$} & \textbf{Simbol PV} & \textbf{Rumus Nilai Sekarang (PV)} & \textbf{Rumus Akumulasi (AV)} \\
\midrule
\textbf{Kontinu Sederhana} & Laju $1$ kontinu & $\bar{a}_{\overline{n}|}$ & $\dfrac{1 - v^n}{\delta}$ & $\dfrac{(1+i)^n - 1}{\delta} = \bar{s}_{\overline{n}|}$ \\[1em]
\textbf{Aritmatika Umum} & $P + (t-1)Q$ & $A$ & $P a_{\overline{n}|} + Q \left(\dfrac{a_{\overline{n}|} - n v^n}{i}\right)$ & $P s_{\overline{n}|} + Q \left(\dfrac{s_{\overline{n}|} - n}{i}\right)$ \\[1em]
\textbf{Meningkat Diskrit (Akhir)} & $1, 2, 3, \dots, n$ & $(Ia)_{\overline{n}|}$ & $\dfrac{\ddot{a}_{\overline{n}|} - n v^n}{i}$ & $\dfrac{\ddot{s}_{\overline{n}|} - n}{i} = (Is)_{\overline{n}|}$ \\[1em]
\textbf{Meningkat Diskrit (Awal)} & $1, 2, 3, \dots, n$ di muka & $(I\ddot{a})_{\overline{n}|}$ & $\dfrac{\ddot{a}_{\overline{n}|} - n v^n}{d}$ & $\dfrac{\ddot{s}_{\overline{n}|} - n}{d} = (I\ddot{s})_{\overline{n}|}$ \\[1em]
\textbf{Menurun Diskrit (Akhir)} & $n, n-1, \dots, 2, 1$ & $(Da)_{\overline{n}|}$ & $\dfrac{n - a_{\overline{n}|}}{i}$ & $\dfrac{n(1+i)^n - s_{\overline{n}|}}{i} = (Ds)_{\overline{n}|}$ \\[1em]
\textbf{Geometrik ($k \ne i$)} & $1, (1+k), \dots, (1+k)^{n-1}$ & $PV_{\text{geom}}$ & $\dfrac{1 - \left(\frac{1+k}{1+i}\right)^n}{i - k}$ & $\dfrac{(1+i)^n - (1+k)^n}{i - k}$ \\[1em]
\textbf{Geometrik ($k = i$)} & $1, (1+i), \dots, (1+i)^{n-1}$ & $PV_{\text{geom}}$ & $n v = \dfrac{n}{1+i}$ & $n(1+i)^{n-1}$ \\[1em]
\textbf{Bervariasi $k > 1$} & $1$ di $k, 2$ di $2k, \dots$ & $A$ & $\dfrac{\frac{a_{\overline{n}|}}{a_{\overline{k}|}} - \frac{n}{k} v^n}{i s_{\overline{k}|}}$ & $A(1+i)^n$ \\[1em]
\textbf{Meningkat Tahunan ($m$-thly)} & $t/m$ tiap sub-periode & $(Ia)_{\overline{n}|}^{(m)}$ & $\dfrac{\ddot{a}_{\overline{n}|} - n v^n}{i^{(m)}}$ & $\dfrac{\ddot{s}_{\overline{n}|} - n}{i^{(m)}} = (Is)_{\overline{n}|}^{(m)}$ \\[1em]
\textbf{Meningkat Tiap Sub-Periode} & $t/m^2$ tiap sub-periode & $(I^{(m)} a)_{\overline{n}|}^{(m)}$ & $\dfrac{\ddot{a}_{\overline{n}|}^{(m)} - n v^n}{i^{(m)}}$ & $\dfrac{\ddot{s}_{\overline{n}|}^{(m)} - n}{i^{(m)}} = (I^{(m)} s)_{\overline{n}|}^{(m)}$ \\[1em]
\textbf{Meningkat Kontinu} & Laju $f(t) = t$ & $(\bar{I}\bar{a})_{\overline{n}|}$ & $\dfrac{\bar{a}_{\overline{n}|} - n v^n}{\delta}$ & $\dfrac{\bar{s}_{\overline{n}|} - n}{\delta} = (\bar{I}\bar{s})_{\overline{n}|}$ \\[1em]
\textbf{Menurun Kontinu} & Laju $f(t) = n - t$ & $(\bar{D}\bar{a})_{\overline{n}|}$ & $\dfrac{n - \bar{a}_{\overline{n}|}}{\delta}$ & $\dfrac{n(1+i)^n - \bar{s}_{\overline{n}|}}{\delta} = (\bar{D}\bar{s})_{\overline{n}|}$ \\
\bottomrule
\end{tabular}
}
\end{center}

---

\section{Latihan Soal \& Pembahasan Mandiri Berstandar Ujian Aktuaria}

\begin{examplebox}[Latihan Soal 1: Anuitas Tangga Bertingkat (\textit{Step-Up Annuity})]
Suatu dana beasiswa memberikan tunjangan pada setiap akhir tahun selama 15 tahun:
\begin{itemize}
    \item Tahun 1 hingga 5: $\$1.000$ per tahun.
    \item Tahun 6 hingga 10: $\$2.000$ per tahun.
    \item Tahun 11 hingga 15: $\$3.000$ per tahun.
\end{itemize}
Jika suku bunga efektif tahunan adalah $6\%$, tentukan nilai sekarang ($PV$) dari program beasiswa tersebut!
\end{examplebox}

\begin{proof}[\textbf{Pembahasan:}]
Aliran dana dapat dipecah menjadi 3 lapis anuitas:
\begin{itemize}
    \item Lapis 1: $\$1.000$ per tahun selama 15 tahun (tahun 1--15).
    \item Lapis 2: Tambahan $\$1.000$ per tahun selama 10 tahun terakhir (tahun 6--15).
    \item Lapis 3: Tambahan $\$1.000$ per tahun selama 5 tahun terakhir (tahun 11--15).
\end{itemize}
\begin{equation*}
    PV = 1000 \, a_{\overline{15}|0.06} + 1000 \, v^5 \, a_{\overline{10}|0.06} + 1000 \, v^{10} \, a_{\overline{5}|0.06}
\end{equation*}
Pada $i = 0.06$:
\begin{align*}
    a_{\overline{15}|} &= \frac{1 - (1.06)^{-15}}{0.06} \approx 9.712249 \\
    v^5 \, a_{\overline{10}|} &= (1.06)^{-5} \left( \frac{1 - (1.06)^{-10}}{0.06} \right) = 0.747258 \times 7.360087 \approx 5.500000 \\
    v^{10} \, a_{\overline{5}|} &= (1.06)^{-10} \left( \frac{1 - (1.06)^{-5}}{0.06} \right) = 0.558395 \times 4.212364 \approx 2.352199
\end{align*}
\begin{equation*}
    PV = 1000 [ 9.712249 + 5.500000 + 2.352199 ] = 1000 \times 17.564448 = \mathbf{\$17.564,45}
\end{equation*}
\end{proof}

\vspace{1.5em}

\begin{examplebox}[Latihan Soal 2: Ujian Aktuaria FM - Hubungan Anuitas Meningkat dan Menurun]
Diketahui dua jenis anuitas 10 tahun dengan suku bunga efektif tahunan $i > 0$:
\begin{itemize}
    \item Anuitas A adalah anuitas meningkat standar $(Ia)_{\overline{10}|} = 32.687$.
    \item Anuitas B adalah anuitas menurun standar $(Da)_{\overline{10}|} = 41.124$.
\end{itemize}
Tentukan suku bunga efektif tahunan $i$!
\end{examplebox}

\begin{proof}[\textbf{Pembahasan:}]
Gunakan Teorema Relasi Komplementer untuk $n = 10$:
\begin{equation*}
    (Ia)_{\overline{10}|} + (Da)_{\overline{10}|} = 11 \, a_{\overline{10}|} \implies 32.687 + 41.124 = 11 \, a_{\overline{10}|}
\end{equation*}
\begin{equation*}
    a_{\overline{10}|} = \frac{73.811}{11} \approx 6.71009
\end{equation*}
Gunakan formula anuitas menurun:
\begin{equation*}
    (Da)_{\overline{10}|} = \frac{10 - a_{\overline{10}|}}{i} \implies 41.124 = \frac{10 - 6.71009}{i} = \frac{3.28991}{i}
\end{equation*}
\begin{equation*}
    i = \frac{3.28991}{41.124} = \mathbf{0.0800} = \mathbf{8.00\%}
\end{equation*}
\end{proof}

\end{document}
